\dictchar{उ}

\bhojentry{उ}{\K{\bu}}{U}{स्वर / सर्वनाम | तद्भव}{\starsThree}%
{हिंदी तथा संस्कृत का पाँचवाँ स्वर वर्ण; वह (दूर की वस्तु या व्यक्ति)।}%
{उ (Vowel U / Pronoun: वह)}%
{Fifth vowel of Indic alphabet; demonstrative pronoun 'that'.}
\bhojexample{उ हमार पुरान खेत बा।}{That is my old farm.}

\bhojentry{उगना}{\K{\bu\bga\bna\bmaa}}{Ugana}{हि. क्रि. | तद्भव}{\starsThree}%
{उगेल, निकलना, प्रकट होना (जैसे सूर्य या पौधे का उगना)।}%
{उगना / निकलल}%
{To rise, sprout, emerge.}
\bhojexample{पूरब में सुरुज उग गइलन।}{The sun rose in the east.}

\bhojentry{उगावल}{\K{\bu\bga\bmaa\bva\bla}}{Ugaawal}{हि. क्रि. | [ठेठ]}{\starsThree}%
{पैदा करना, उपजाना; कर्ज़ या चंदा वसूल करना।}%
{उगावल / वसूल कइल}%
{To grow, cultivate; to collect dues.}

\bhojsubentry{उगाहने}{\K{\bu\bga\bmaa\bha\bna\bme}}{Ugaahane}{हि. पुं. | [ठेठ]}{\starsTwo}%
{लगान, चंदा या कर्ज़ वसूलने की क्रिया।}%
{उगाहने / वसूली}%
{Collection of tax/dues.}

\bhojentry{उघार}{\K{\bu\bgha\bmaa\bra}}{Ughaar}{हि. वि. | [ठेठ]}{\starsThree}%
{खुला हुआ, नग्न, बिना ढका।}%
{उघार / खुला}%
{Open, bare, uncovered.}

\bhojsubentry{उघारे}{\K{\bu\bgha\bmaa\bra\bme}}{Ughaare}{हि. अव्य. | [ठेठ]}{\starsThree}%
{नंगे, बिना ढके अवस्था में (जैसे उघारे गोड़)।}%
{उघारे (Ughaare) / नंगे}%
{Barefoot, bare, exposed.}
\bhojexample{उघारे गोड़ घाम में मत चला।}{Do not walk barefoot in the sun.}

\bhojsubentry{उघारल}{\K{\bu\bgha\bmaa\bra\bla}}{Ughaaral}{हि. क्रि. | [ठेठ]}{\starsThree}%
{ढक्कन या पर्दा हटाना, खोलना।}%
{उघारल / खोलल}%
{To uncover, reveal, open.}

\bhojentry{उचकना}{\K{\bu\bca\bka\bna\bmaa}}{Uchakana}{हि. क्रि. | तद्भव}{\starsTwo}%
{एड़ी उठाकर ऊपर देखना या खड़े होना।}%
{उचकना}%
{To stand on tiptoe, jump slightly.}

\bhojsubentry{उचकावल}{\K{\bu\bca\bka\bmaa\bva\bla}}{Uchakaawal}{हि. क्रि. | [ठेठ]}{\starsTwo}%
{ऊपर उठाना, उचकाना।}%
{उचकावल}%
{To lift up, raise slightly.}

\bhojentry{उछाह}{\K{\bu\bcha\bmaa\bha}}{Uchhaah}{हि. पुं. | तद्भव}{\starsThree}%
{उत्साह, हर्ष, उमंग, उल्लास।}%
{उछाह (Uchhaah) / उमंग}%
{Enthusiasm, joy, excitement.}
\bhojexample{शादी के घर में भारी उछाह बा।}{There is great joy in the wedding house.}

\bhojentry{उजाड़}{\K{\bu\bja\bmaa\bdda\bnukta}}{Ujaar}{हि. वि./पुं. | तद्भव}{\starsThree}%
{विरान, सुनसान, नष्ट-भ्रष्ट स्थान।}%
{उजाड़ / विरान}%
{Desolate, ruined, abandoned place.}

\bhojsubentry{उजरावल}{\K{\bu\bja\bra\bmaa\bva\bla}}{Ujaraawal}{हि. क्रि. | [ठेठ]}{\starsThree}%
{नष्ट करना, बर्बाद करना, उजाड़ना।}%
{उजरावल / बर्बाद कइल}%
{To devastate, ruin, lay waste.}

\bhojentry{उजला}{\K{\bu\bja\bla\bmaa}}{Ujala}{हि. वि. | तद्भव}{\starsThree}%
{सफ़ेद, शुभ्र, चमकदार।}%
{उजला / सफ़ेद}%
{White, bright, clean.}

\bhojsubentry{उजियार}{\K{\bu\bja\bmi\bya\bmaa\bra}}{Ujiyaar}{हि. पुं. | [ठेठ]}{\starsThree}%
{रोशनी, प्रकाश, उजाला।}%
{उजियार / उजाला}%
{Light, brightness.}
\bhojsee{इजोत}

\bhojentry{उतरना}{\K{\bu\bta\bra\bna\bmaa}}{Utarana}{हि. क्रि. | तद्भव}{\starsThree}%
{नीचे आना, उतरल (जैसे सवारी से उतरना)।}%
{उतरना / नीचे आइल}%
{To descend, get down, dismount.}

\bhojsubentry{उतराव}{\K{\bu\bta\bra\bmaa\bva}}{Utaraav}{हि. पुं. | तद्भव}{\starsTwo}%
{ढलान, नीचे उतरने का रास्ता।}%
{उतराव / ढलान}%
{Slope, descent.}

\bhojsubentry{उतार}{\K{\bu\bta\bmaa\bra}}{Utaar}{हि. पुं. | तद्भव}{\starsThree}%
{उतारने का काम; ढाल; उतरा हुआ कपड़ा।}%
{उतार / ढलान}%
{Decline, slope, reduction.}

\bhojsubentry{उतारू}{\K{\bu\bta\bmaa\bra\bmuu}}{Utaaru}{हि. वि. | तद्भव}{\starsThree}%
{आमादा, तत्पर (जैसे झगड़ा पर उतारू)।}%
{उतारू / तैयार}%
{Bent upon, determined, eager.}

\bhojentry{उतावला}{\K{\bu\bta\bmaa\bva\bla\bmaa}}{Utaawala}{हि. वि. | तद्भव}{\starsThree}%
{व्यग्र, आतुर, हड़बड़िया।}%
{उतावला / हड़बड़िया}%
{Impatient, hasty, restless.}

\bhojsubentry{उतावली}{\K{\bu\bta\bmaa\bva\bla\bmii}}{Utaawali}{हि. स्त्री. | तद्भव}{\starsThree}%
{हड़बड़ी, जल्दबाजी।}%
{उतावली / जल्दबाजी}%
{Haste, impatience.}

\bhojentry{उतपत}{\K{\bu\bta\bpa\bta}}{Utapat}{हि. पुं. | [ठेठ]}{\starsTwo}%
{उत्पात, उपद्रव, हुड़दंग।}%
{उत्पात / खुराफात}%
{Nuisance, commotion, mischief.}

\bhojentry{उत्साह}{\K{\bu\bta\bvir\bsa\bmaa\bha}}{Utsaah}{सं. पुं. | तत्सम}{\starsThree}%
{जोश, उछाह, उमंग।}%
{जोश / उछाह}%
{Zeal, enthusiasm, spirit.}

\bhojentry{उद्धार}{\K{\bu\bdha\bvir\bdha\bmaa\bra}}{Uddhaar}{सं. पुं. | तत्सम}{\starsThree}%
{मुक्ति, संकट से छुटकारा, कल्याण।}%
{उद्धार / मुक्ति}%
{Salvation, rescue, deliverance.}

\bhojentry{उधार}{\K{\bu\bdha\bmaa\bra}}{Udhaar}{हि. पुं. | तद्भव}{\starsThree}%
{कर्ज़, ऋण, पैंचा।}%
{उधार / कर्ज़}%
{Credit, loan, debt.}
\bhojexample{दुकानदार से उधार पर सामान मिल गइल।}{Goods were obtained on credit from the shopkeeper.}

\bhojsubentry{उधारी}{\K{\bu\bdha\bmaa\bra\bmii}}{Udhaari}{हि. स्त्री. | तद्भव}{\starsThree}%
{बाकी कर्ज़, उधारी का लेन-देन।}%
{उधारी}%
{Credit transaction, outstanding debt.}

\bhojentry{उपज}{\K{\bu\bpa\bja}}{Upaj}{हि. स्त्री. | तद्भव}{\starsThree}%
{फसल की पैदावार, उत्पादन।}%
{उपज / पैदावार}%
{Crop yield, produce, harvest.}

\bhojsubentry{उपजाऊ}{\K{\bu\bpa\bja\bmaa\bmuu}}{Upajaau}{हि. वि. | तद्भव}{\starsThree}%
{उर्वर, उपजाऊ खेत।}%
{उपजाऊ / उर्वर}%
{Fertile, productive.}

\bhojentry{उपाय}{\K{\bu\bpa\bmaa\bya}}{Upaay}{सं. पुं. | तत्सम}{\starsThree}%
{जुगत, तरकीब, युक्ति।}%
{जुगत (Jugat) / तरकीब}%
{Remedy, solution, means.}

\bhojentry{उपहार}{\K{\bu\bpa\bha\bmaa\bra}}{Upahaar}{सं. पुं. | तत्सम}{\starsThree}%
{भेंट, सौगात, पुरस्कार।}%
{भेंट / सौगात}%
{Gift, present.}

\bhojentry{उबटन}{\K{\bu\bba\btta\bna}}{Ubatan}{हि. पुं. | तद्भव}{\starsThree}%
{सरसों-हल्दी की पीठी, बुकवा।}%
{बुकवा (Bukwa) / पीठी}%
{Herbal cosmetic paste, ubtan.}

\bhojentry{उबाल}{\K{\bu\bba\bmaa\bla}}{Ubaal}{हि. पुं. | तद्भव}{\starsThree}%
{खौलना, फान, उबाल।}%
{उबाल / फान}%
{Boiling, ebullition.}

\bhojsubentry{उबालना}{\K{\bu\bba\bmaa\bla\bna\bmaa}}{Ubaalana}{हि. क्रि. | तद्भव}{\starsThree}%
{पानी में खौलाना, उसिना।}%
{उबालना / उसिना}%
{To boil, parboil.}

\bhojentry{उबारल}{\K{\bu\bba\bmaa\bra\bla}}{Ubaaral}{हि. क्रि. | [ठेठ]}{\starsTwo}%
{संकट से उबारना, रक्षा करना।}%
{उबारल / बचावल}%
{To rescue, liberate, save.}

\bhojentry{उमंग}{\K{\bu\bma\banus\bga}}{Umang}{हि. स्त्री. | तद्भव}{\starsThree}%
{उत्साह, मौज, तरंग।}%
{उमंग / मौज}%
{Joyous excitement, wave of joy.}

\bhojentry{उमर}{\K{\bu\bma\bra}}{Umar}{हि. स्त्री. | अरबी आगत}{\starsThree}%
{आयु, अवस्था, उमरिया।}%
{उमर / आयु}%
{Age, lifespan.}

\bhojsubentry{उमरिया}{\K{\bu\bma\bra\bmi\bya\bmaa}}{Umariya}{हि. स्त्री. | [ठेठ]}{\starsThree}%
{उमर, आयु (जैसे छोटी उमरिया)।}%
{उमरिया / उमर}%
{Age, young age.}
\bhojexample{छोटी उमरिया में लइका बहुत चतुर बा।}{At a young age, the boy is very clever.}

\bhojentry{उमड़ना}{\K{\bu\bma\bdda\bnukta\bna\bmaa}}{Umarana}{हि. क्रि. | तद्भव}{\starsThree}%
{घिरना, उमड़ना (जैसे बादल उमड़ल)।}%
{उमड़ल / घिरल}%
{To gather, surge, swell (as clouds).}

\bhojentry{उमस}{\K{\bu\bma\bsa}}{Umas}{हि. स्त्री. | तद्भव}{\starsThree}%
{गर्मी की चिपचिपाहट, घमस।}%
{उमस / घमस}%
{Sultriness, humid heat.}

\bhojentry{उरहन}{\K{\bu\bra\bha\bna}}{Urahan}{हि. पुं. | [ठेठ]}{\starsThree}%
{उलाहना, शिकायत, सिकायत।}%
{उरहन (Urahan) / शिकायत}%
{Complaint, reproach.}
\bhojexample{पड़ोसी घरे उरहन लेके आइलन।}{The neighbor came home with a complaint.}

\bhojentry{उलझन}{\K{\bu\bla\bjha\bna}}{Ulahan}{हि. स्त्री. | तद्भव}{\starsThree}%
{उधेड़बुन, फँसाव, पेच, असमंजस।}%
{उलझन / पेच}%
{Entanglement, confusion, problem.}

\bhojsubentry{उलझावल}{\K{\bu\bla\bjha\bmaa\bva\bla}}{Ulahaawal}{हि. क्रि. | [ठेठ]}{\starsThree}%
{फँसाना, पेचीदा करना।}%
{उलझावल}%
{To entangle, complicate.}

\bhojentry{उलटा}{\K{\bu\bla\btta\bmaa}}{Ulata}{हि. वि. | तद्भव}{\starsThree}%
{औंधा, विपरीत, पुट (जैसे कपड़ा उलटा बा)।}%
{उलटा / पुट}%
{Inverted, reverse, opposite.}

\bhojsubentry{उलटना}{\K{\bu\bla\btta\bna\bmaa}}{Ulatana}{हि. क्रि. | तद्भव}{\starsThree}%
{पलटना, औंधाना।}%
{उलटना / पलटल}%
{To overturn, turn upside down.}

\bhojentry{उलू}{\K{\bu\bla\bmuu}}{Uloo}{हि. पुं. | तद्भव}{\starsThree}%
{उल्लू पक्षी; निहायत मूर्ख व्यक्ति।}%
{उल्लू (Owl) / मूरुख}%
{Owl; fool.}

\bhojentry{उल्लास}{\K{\bu\bla\bvir\bla\bmaa\bsa}}{Ullaas}{सं. पुं. | तत्सम}{\starsThree}%
{आनंद, उछाह, ख़ुशी।}%
{आनंद / उछाह}%
{Joy, delight, exultation.}

\bhojentry{उसाँस}{\K{\bu\bsa\bmaa\banus\bsa}}{Usaañs}{हि. स्त्री. | [ठेठ]}{\starsThree}%
{ठंडी सांस, आह, उसास।}%
{उसाँस (Usaañs) / आह}%
{Sigh, deep breath.}

\bhojentry{उसावल}{\K{\bu\bsa\bmaa\bva\bla}}{Usaawal}{हि. क्रि. | [ठेठ/कृषि]}{\starsThree}%
{हवा में अनाज ऊपर से गिराकर भूसा अलग करना।}%
{उसावल (Usaawal)}%
{To winnow grain in wind.}
\bhojexample{खलिहान में किसान गूँह उसावत बाड़न।}{In the barn, the farmer is winnowing wheat.}

\bhojentry{उसिना चावल}{\K{\bu\bsa\bmi\bna\bmaa} \ \K{\bca\bmaa\bva\bla}}{Usina chaawal}{हि. पुं. | [ठेठ]}{\starsThree}%
{उबला हुआ पक्का धान का चावल (उसिना)।}%
{उसिना चावल (Parboiled Rice)}%
{Parboiled rice.}

\bhojentry{उस्ताद}{\K{\bu\bsa\bvir\bta\bmaa\bda}}{Ustaad}{हि. पुं. | फारसी आगत}{\starsThree}%
{गुरु, शिक्षक, कला में माहिर।}%
{गुरु / उस्ताद}%
{Master, expert, teacher.}

\bhojentry{उस्तरा}{\K{\bu\bsa\bvir\bta\bra\bmaa}}{Ustara}{हि. पुं. | फारसी आगत}{\starsThree}%
{हजामत बनाने का तीखा छुरा।}%
{उस्तरा}%
{Razor.}

\bhojentry{उहवाँ}{\K{\bu\bha\bva\banus}}{Uhawañ}{हि. अव्य. | [ठेठ]}{\starsThree}%
{वहाँ, उस स्थान पर।}%
{उहवाँ / वहाँ}%
{There, in that place.}

\bhojentry{उहे}{\K{\bu\bha\bme}}{Uhe}{हि. सर्व. | [ठेठ]}{\starsThree}%
{वही (बात, स्थिति या कथन)।}%
{उहे / वही}%
{That very thing/statement.}

\bhojentry{उही}{\K{\bu\bha\bmii}}{Uhi}{हि. सर्व. | [ठेठ]}{\starsThree}%
{वही विशिष्ट व्यक्ति या प्राणी।}%
{उही / वही}%
{That specific person right there.}

\bhojentry{उहो}{\K{\bu\bha\bmo}}{Uho}{हि. सर्व. | [ठेठ]}{\starsThree}%
{वह भी।}%
{उहो / वह भी}%
{That as well.}

\bhojentry{उ-सभ}{{\K{\bu}}-\K{\bsa\bbha}}{U-sabh}{हि. सर्व. | [ठेठ]}{\starsThree}%
{वे सब (बहुवचन)।}%
{उ-सभ / वे सब}%
{All those.}

\bhojentry{उ-ठवाँ}{{\K{\bu}}-\K{\bttha\bva\banus}}{U-thavan}{हि. अव्य. | [ठेठ]}{\starsThree}%
{उस स्थान पर, उसी जगह।}%
{उ-ठवाँ / ओह जगह}%
{In that exact place.}
